% Source: Markdown/close_vs_open_loop.md
\begin{notechapter}
\chapter{Open-Loop versus Closed-Loop Prediction}\label{ch:open-closed-loop}

In time series prediction, \textbf{open-loop} and \textbf{closed-loop}
refer to two distinct approaches for generating forecasts, differing in
whether \textbf{true observed values} are used to refine predictions
during the forecasting process.

\begin{remark}[Two opposite conventions]
The convention below follows control engineering, where a closed loop feeds \emph{measured} outputs
back. In the neural-network time-series literature, in particular for NARX networks (e.g.\ MATLAB's
\texttt{narxnet}), the words are used the other way round: the \emph{open-loop} (series--parallel)
configuration feeds the \emph{true} past outputs to the network (used for training and one-step-ahead
prediction), while the \emph{closed-loop} (parallel) configuration feeds the network's \emph{own
predictions} back as inputs (used for multi-step forecasting). Always check which convention a source
uses; what matters is whether the inputs at each step are observed values or the model's own
predictions.
\end{remark}

\section{Open-Loop (Feedforward) Prediction}

\begin{itemize}
\tightlist
\item
  \textbf{Definition}: Predictions are generated in a single pass, using
  only the \textbf{initial input sequence}. Future steps are forecasted
  based on the model's own previous predictions, without incorporating
  actual observations once prediction begins.
\item
  \textbf{Process}:

  \begin{itemize}
  \tightlist
  \item
    The model predicts multiple future steps (e.g.,
    \(t+1, t+2, ..., t+n\)) at once.
  \item
    Each step \(t+k\) relies on the predicted value of \(t+(k-1)\), not
    the true observation.
  \end{itemize}
\item
  \textbf{Pros}: Efficient for batch predictions; suitable when
  immediate feedback is unavailable.
\item
  \textbf{Cons}: \textbf{Error accumulation} due to reliance on prior
  predictions. Errors propagate as the forecast horizon increases.
\item
  \textbf{Example}: Predicting the next 10 days of stock prices all at
  once, using only historical data up to today.
\end{itemize}

\begin{center}\rule{0.5\linewidth}{0.5pt}\end{center}

\section{Closed-Loop (Feedback) Prediction}

\begin{itemize}
\tightlist
\item
  \textbf{Definition}: Predictions are updated iteratively by
  incorporating \textbf{true observed values} as they become available.
  Each new forecast uses the latest ground-truth data to ``correct'' the
  model.
\item
  \textbf{Process}:

  \begin{itemize}
  \tightlist
  \item
    Predict one step ahead (e.g., \(t+1\)).
  \item
    Once the true value at \(t+1\) is observed, it is fed back into the
    model to predict \(t+2\), and so on.
  \end{itemize}
\item
  \textbf{Pros}: Reduces error propagation by resetting the input with
  actual data at each step.
\item
  \textbf{Cons}: Requires real-time access to ground-truth observations,
  which may not always be feasible.
\item
  \textbf{Example}: Weather forecasting updated daily with the latest
  measurements.
\end{itemize}

\begin{center}\rule{0.5\linewidth}{0.5pt}\end{center}

\section{Key Difference}

\begin{longtable}[]{@{}
  >{\raggedright\arraybackslash}p{(\columnwidth - 4\tabcolsep) * \real{0.2407}}
  >{\raggedright\arraybackslash}p{(\columnwidth - 4\tabcolsep) * \real{0.3704}}
  >{\raggedright\arraybackslash}p{(\columnwidth - 4\tabcolsep) * \real{0.3889}}@{}}
\toprule\noalign{}
\begin{minipage}[b]{\linewidth}\raggedright
\textbf{Aspect}
\end{minipage} & \begin{minipage}[b]{\linewidth}\raggedright
\textbf{Open-Loop}
\end{minipage} & \begin{minipage}[b]{\linewidth}\raggedright
\textbf{Closed-Loop}
\end{minipage} \\
\midrule\noalign{}
\endhead
\bottomrule\noalign{}
\endlastfoot
\textbf{Feedback} & No use of true observations post-start & Uses true
observations at each step \\
\textbf{Error Propagation} & High (compounding errors) & Low (errors
reset with true data) \\
\textbf{Use Case} & Batch forecasting & Real-time/online forecasting \\
\end{longtable}

\begin{center}\rule{0.5\linewidth}{0.5pt}\end{center}

\section{Practical Applications}

\begin{itemize}
\tightlist
\item
  \textbf{Open-Loop}: Long-term scenario planning, energy demand
  forecasting, or any context where immediate feedback is impractical.
\item
  \textbf{Closed-Loop}: Autonomous systems (e.g., robotics), real-time
  financial trading, or adaptive control systems.
\end{itemize}

Understanding this distinction helps in selecting the right approach
based on data availability and the need for accuracy versus
computational efficiency.
\end{notechapter}
